% RQ2 Conclusion: Governance Capture Mitigation

\subsection{Summary}

We have presented a systematic comparison of governance capture mitigation strategies through multi-agent simulation. Across \experimentruns{} simulation runs testing 27 configurations of vote caps, quadratic voting, and velocity penalties, we characterized the effectiveness and tradeoffs of each approach.

Key findings:

\begin{enumerate}
    \item \textbf{Vote caps}: Effective at limiting peak influence but create sybil incentives. Best applied to delegation rather than direct holdings.

    \item \textbf{Quadratic voting}: Provides smooth power compression, roughly halving the Gini coefficient. Vulnerable to sybil without identity systems.

    \item \textbf{Velocity penalties}: Effective against rapid accumulation attacks; 30-60 day windows balance protection with inclusion.

    \item \textbf{Hybrid approaches}: Combinations outperform single mechanisms. Recommended: quadratic base + delegation caps + velocity buffer.

    \item \textbf{Tradeoff frontier}: All mitigation imposes some governance efficiency cost. Optimal configuration depends on DAO priorities.
\end{enumerate}

\subsection{Contributions}

This paper contributes:

\begin{enumerate}
    \item \textbf{Quantified comparison} of capture mitigation mechanisms under controlled conditions
    \item \textbf{Formal metrics} for capture risk assessment (whale influence, Nakamoto coefficient, composite score)
    \item \textbf{Parameter guidelines} for each mechanism based on simulation evidence
    \item \textbf{Hybrid recommendations} for layered capture defense
\end{enumerate}

\subsection{Implications}

For DAO designers, our results provide evidence-based guidance for mechanism selection. Pure token voting is vulnerable to capture; mitigation mechanisms are not optional for governance resilience.

For the research community, our framework enables continued investigation of capture dynamics. The interaction between mitigation mechanisms and adversarial behavior deserves ongoing attention.

For the ecosystem, our findings underscore that decentralization is not automatic. Token distribution creates power distribution; mechanism design determines whether that power is checked.

\subsection{Closing Remarks}

Governance capture is not theoretical---it is observable in deployed DAOs where small groups dominate decision-making. The mechanisms we evaluate are not exotic proposals but practical tools available today.

The choice is not whether to implement capture mitigation, but which combination of mechanisms best serves the DAO's values and operational needs. Our simulations illuminate these tradeoffs, providing a foundation for informed governance design.

We release our simulation framework and experimental configurations to enable the community to extend this analysis, test additional mechanisms, and develop the evidence base for capture-resistant governance.
